\section{Characters sums of \texorpdfstring{$\psi(f(X))\chi(g(X))$}{Ψ(f(X))χ(g(X))}}
In this section we will use the multiplicative function $\xi:G\to \bbC$ in the previous section and prove bounds on character sums like $\sum_{\alpha\in\bbF_{q^k}}\psi^{(k)}(f(X))\chi^{(k)}(g(X))$ where $\chi\in\Xq$ and $\psi\in\Mq$ and they are lifted to $\chi^{(k)}$ and $\psi^{(k)}$ for the field $\bbF_{q^k}$. 

So suppose we have a additive character $\chi\in\Xq$ and a multiplicative character $\psi\in\Mq$ of $\bbF_q$. Let $f(X)\in\bbF_q[X]$ be a monic polynomial in $\bbF_q[X]$. Let $f(X)$ factors into $$f(X)=(X-\gm_1)^{a_1}\cdots (X-\gm_m)^{a_m}$$ in the field $\ov{\bbF}_q[X]$ i.e. $f$ has $m$ distinct roots. And let $g\in\bbF_q[X]$ be any fixed polynomial of degree $n$ and constant term is $0$. Define the $G, \ovG, \tdH, H$ and $\ovH$ like previous section. Then we define the maps $\vsg:G\to \bbF_q$ and $\vrh:G\to\bbF_q$ and the corresponding multiplicative map $\xi:G\to\bbC$.
\subsection{\texorpdfstring{$L$}{L}-function on \texorpdfstring{$\xi$}{ξ}}
We want to create the $L$-function on $\xi$ and use to calculate the character sums and get bounds on absolute value of the character sum. We will show that the $L$-function $\ovL(z,\xi)$ is of finite degree. First we will show $\xi$ is indeed non-trivial. 
\begin{Theorem}{Non-triviality of $\xi$}{thm:xi-nontrivial}
  Suppose either $\psi\in\Mq$ is a nontrivial multiplicative character of exponent $d$ and $Y^d-f(X)$ is absolutely irreducible. If the map $\xi$ is non-trivial i.e. $\exs\ r\in \ovG$ such that $\chi(r)\neq 1$. 
\end{Theorem}
\begin{proof}
  Suppose $\xi(r)=1$ for all $r\in\ovG$. Then $(\psi\circ\vsg)(r)\cdot (\chi\circ \vrh)(r)=1$. Since $\psi\circ\vsg(r)$ is a $d^{th}$ root of unity and $\chi\circ \vrh(r)$ is a $p^{th}$ roots of unity with $\gcd(d,p)=1$ we have $$\psi\circ\vsg(r)=\chi\circ\vrh(r)=1$$
  Therefore it suffices to find a $k\in \ovG$ such that $\psi\circ\vsg(k)\neq 1$ in the first case and in the second case it will suffice to find a $k\in\ovG$ such that $\chi\circ\vrh(k)\neq 1$. 

  \paragraph*{$\psi$ is non-trivial:} Suppose $\ord(\psi)=e$ where $e\mid d$. Since $Y^d-f(X)$ is absolutely irreducible not all $a_i$ are multiplies of $e$ by \thmref{thm:yd-fx-irr}. Without loss of generality suppose $e\nmid a_1$. For any $c_2,\dots, c_m\in\bbF_q^*$ we can pick $c_1\in\bbF_q^*$ such that $c_1^{a_1}\cdots c_m^{a_m}\notin \lt(\bbF_q^*\rt)^{(e)}$ and hence $\psi\lt(c_1^{a_1}\cdots c_m^{a_m}\rt)\neq 1$. Now by the arguments of \lemref{lem:ql-poly-every-coset} there exists a polynomial $k(X)\in\ovG$ such that $k(\gm_i)=c_i$ for all $i\in[m]$. Therefore $\vsg(k)=c_1^{a_1}\cdots c_m^{a_m}$ and hence $\psi\circ \vsg(k)\neq 1$.
  \paragraph*{$\chi$ is non-trivial:} \begin{enumerate}[label=(\roman*)]
    \item Suppose first that $\gcd(n,q)=1$. Then suppose $g(X)=aX^n+g_1(X)$ where $\deg(g_1)<n$ and $g_1(0)=0$. If $k(X)=X^n+\alpha=\prod_{j=1}^n(X-\alpha_j)$ where $\alpha_j\in\ov{\bbF_q}$. Then $g_1(\alpha_1)+\cdots g_n(\alpha_n)=0$ since $g_1$ is a polynomial with constant term $0$ and it is a polynomial in the first $n-1$ elementary symmetric polynomials in $\alpha_1,\dots, \alpha_n$. Therefore $$\sum_{j=1}^ng(\alpha_j)=a\sum_{j=1}^n\alpha_j^n=a\cdot n\cdot (-1)^{n+1}\cdot \alpha$$Therefore $\vrh(k)=a\cdot n\cdot (-1)^{n+1}\cdot \alpha$ and thus $\chi\circ\vrh(k)=\chi\lt(a\cdot n\cdot (-1)^{n+1}\cdot \alpha\rt)$. So for a proper choice of $\alpha$ we get $\chi\circ\vrh(k)\neq 1$ since $n$ is not divisible by the $\Char(\bbF_q)$.
    \item Suppose $Y^q-Y-f(X)$ is absolutely irreducible. Now for any $\alpha\in\bbF_q^*$ we have $$\alpha\cdot Y^q-\alpha\cdot Y-g(X)=(\alpha\cdot Y)^q-(\alpha\cdot Y)-f(X)$$ therefore $\alpha\cdot Y^q-\alpha\cdot Y-g(X)$ is absolutely irreducible. Hence $Y^p-Y-\alpha\cdot g(X)$ is also absolutely irreducible where $p=\Char(\bbF_q)$. Let $a$ is such that $\chi(\gm)=\chi_a(\gm)$ for all $\gm\in\bbF_q$ as defined in \AddCharDefthm. Let $Z_r(Y^q-Y-a\cdot g(X))$ is the number of solutions of $Y^q-Y-a\cdot g(X)$ over $\bbF_{q^r}$. Now by \thmref{thm:bombieri} we have $Z_r(Y^q-Y=g(X))=q^r+O(q^{\nicefrac{r}2})$. Hence if $r$ is large $Z_r(Y^q-Y=g(X))<p\cdot q^r$. Let $E$ denote the field $\bbF_{q^r}$. Now suppose $c\in \bbF_{q^r}$. If $\tr_{E}(a\cdot g(c))=0$ then there are $p$ values of $b\in $ such that $b^q-b-a\cdot g(c)=0$. If $\tr_E(a\cdot g(c))\neq 0$ there are no such $b\in\bbF_{q^r}$. Since $Z_r(Y^q-Y=g(X))<p\cdot q^r$ there is a $b\in\bbF_{q^r}$ with $\tr_E(a\cdot g(c))\neq 0$. So take $k(X)=\prod_{i=0}^{r-1}(X-c^{q^i})$. Then $$\vrh(k)=g(c)+g(c^q)+\cdots g(c^{q^{r-1}})=\tr_E(g(c))$$So $\chi\circ \vrh(k)\neq 1$. 
  \end{enumerate} 
  Therefore we have the theorem. 
\end{proof}
We eventually want to construct the $L$-function using the multiplicative function $\xi$. The above theorem helps us to invoke \LFactthm. 
\begin{lemma}{Vanishing of Character Sum on $\xi$}{lem:xi-sum-vanishes}
  Suppose either $\psi\in\Mq$ is a nontrivial multiplicative character of exponent $d$ and $Y^d-f(X)$ is absolutely irreducible. Or, $\chi\in\Xq$ is a non-trivial additive character and
  \begin{enumerate}[label=(\roman*)]
    \item either, $\gcd(n,q)=1$
    \item or, more generally $Y^q-Y-g(X)$ is absolutely irreducible.
  \end{enumerate}
  Suppose $t\geq 0$. Then $$\sum_{\substack{h\in \Phi\\ \deg(h)=n+m+t}}\xi(g)=0$$
\end{lemma}
\begin{proof}
  By \corref{cor:xi-trivial-subgroup} and \thmref{thm:xi-nontrivial}, $\xi$ induces a non-trivial character on the finite group $\quotient{\ovG}{H}$. As $h$ runs through all polynomials of $\ovG$ of degree $n+m+t$ then by \lemref{lem:ql-poly-every-coset} it will lie precisely land $q^l$ times in every coset of $\quotient{\ovG}{H}$. So by \AddCharPropsthm[(i)] and \MultCharPropsthm[(i)] we have the lemma. 
\end{proof}

As we mentioned in the previous section. We can extend $\xi$ to $G$ by setting $\xi(h)=0$ if $h$ is a polynomial not in $\ovG$. So now we form the $L$-functions, $L(s,\xi)$ and $\ovL(Z,\xi)$ on $\xi$. And using the above lemma we have the following theorem:
\begin{Theorem}{Finite Degree of $\ovL(z,\xi)$}{thm:ovL-xi-finite}
  Suppose either $\psi\in\Mq$ is a nontrivial multiplicative character of exponent $d$ and $Y^d-f(X)$ is absolutely irreducible. Or, $\chi\in\Xq$ is a non-trivial additive character and
  \begin{enumerate}[label=(\roman*)]
    \item either, $\gcd(n,q)=1$
    \item or, more generally $Y^q-Y-g(X)$ is absolutely irreducible.
  \end{enumerate}
  Then $$\ovL(z,\xi)=1+c_1z+\cdots+c_{n+m-1}z^{n+m-1},\quad L(s,\xi)=1+c_1u+\cdots+c_{n+m-1}u^{n+m-1}$$where $u=q^{-s}$.
\end{Theorem}

So we can use \LFactthm and factorize $\ovL(z,\xi)$ and get complex numbers $w_1,\dots, w_{n+m-1}$ such that $$\ovL(z,\xi)=\prod_{j=1}^{n+m-1}(1-w_iz).$$ We can say a lot about the coefficients, specifically $c_1$. 
\begin{Theorem}{Leading Coefficient $c_1$ of $\ovL(z,\xi)$}{thm:ovL-leading-coeff}
  If $\psi$ is nontrivial or if $\psi$ is trivial and $f(X)=1$ then $$c_1=\sum_{\gm\in\bbF_q}\psi(f(\gm))\chi(g(\gm)).$$
\end{Theorem}
\begin{proof}
  We have \begin{align*}
    c_1& = \sum_{\substack{h\in\Phi+1\\ \deg(g)=t}}\xi(h)\\ 
    & = \sum_{\substack{\alpha\in\bbF_q\\ \alpha\neq \gm_i\ \forall i\in[m]}}\xi(X-\alpha)\\ 
    & = \sum_{\substack{\alpha\in\bbF_q\\ \alpha\neq \gm_i\ \forall i\in[m]}}\psi\circ\vsg(X-\alpha)\cdot \chi\circ\vrh(X-\alpha)\\ 
    & = \sum_{\substack{\alpha\in\bbF_q\\ f(\alpha)\neq 0}}\psi(f(\alpha))\cdot \chi(g(\alpha))\\ 
    & = \sum_{\alpha\in\bbF_q}\psi(f(\alpha))\cdot \chi(g(\alpha))
  \end{align*}
  So we have the theorem. 
\end{proof}

\subsection{Field Extensions}
Suppose we aer given a finite field extension $E$ of $\bbF_q$ where $[E\colon \bbF_q]=r$. So like \autoref{sec:lifting-chars} we can lift $\psi$ and $\chi$ to $\psi^{(r)}$ and $\chi^{(r)}$ for $E$. And using these two characters we can define the lifted multiplicative function $\xi^{(r)}$ on the group $G_r,\ovG_r$ with rational functions defined similarly as $G,\ovG$ but on $E$. So for any $r\in G_r$ we define $$\xi^{(r)}=\psi^{(r)}\circ \vsg(r)\cdot \chi^{(r)}\circ\vrh(r)$$We define $\tdH_r,H_r,\ovH_r$ for $G_r,\ovG_r$ similarly as $\tdH,H,\ovH$. Now using $\xi^{(r)}$ we also create the $L$-functions, $L^{(r)}(s,\xi^{(r)})$ and $\ovL^{(r)}(z,\xi^{(r)})$. The norm of a polynomial $h\in\Phi^{(r)}$ over $E$ is $\mfn^{(r)}(h)\coloneqq q^{r\cdot\deg(h)}=|E|^{\deg(h)}$, the analogue of $\mfn(h)=q^{\deg(h)}$ for the base field. So we have $$L^{(r)}(s,\xi^{(r)})=\sum_{h\in\Phi^{(r)}}\xi^{(r)}(h)\cdot\mfn^{(r)}(h)^{-s},\qquad \ovL^{(r)}(z,\xi^{(r)})=\sum_{h\in\Phi^{(r)}}\xi^{(r)}(h)\cdot z^{\deg(h)}$$ where $\Phi^{(r)}$ is the set of all monic polynomials over $E$. We similarly define $\Phi_d$ for any $d\in\bbZ_0$. We also use the Weil sum notation from \autoref{sec:character-sum-poly}: $$W(\psi,\chi;f,g)\coloneqq\sum_{\alpha\in\bbF_q}\psi(f(\alpha))\cdot \chi(g(\alpha)),\qquad W^{(r)}(\psi^{(r)},\chi^{(r)};f,g)\coloneqq\sum_{\alpha\in E}\psi^{(r)}(f(\alpha))\cdot\chi^{(r)}(g(\alpha))$$where the superscript $(r)$ indicates the sum is over the extension field $E=\bbF_{q^r}$ using the lifted characters.
\begin{lemma}{Splitting of Irreducible Polynomials over $\bbF_{q^k}$}{lem:irr-splits-over-ext}
  Let $K=\bbF_{q^k}$ and $E=\bbF_{q^r}$. Suppose $h(X)=X^d-a_1X^{d-1}+\cdots+(-1)^da_d$ be an irreducible polynomial in $\bbF_q[X]$. Then in $\bbF_{q^k}[X]$, $h$ splits into $r=\gcd(k,d)$ many irreducible polynomials of degree $\nicefrac{d}{r}$: $$h(X)=h_1(X)\cdots h_r(X)$$If we normalize $h_i(X)$ such that $h_i$ is monic then $h_i(X)\in\bbF_{q^r}[X]$ for all $i\in[r]$. Then
  \begin{enumerate}[label=(\roman*)]
    \item $\mfn^{(k)}(h_i)=\mfn(h)^{\nicefrac{k}{r}}$
    \item $\xi^{(k)}(h_i)=\xi\st^{\nicefrac{k}r}(h)$
  \end{enumerate}
\end{lemma}
\begin{proof}
  \begin{enumerate}[label=(\roman*)]
    \item Now by definition $\mfn^{(k)}(h_i)=q^{k\lt(\nicefrac{d}{r}\rt)}=\mfn(h)^{\nicefrac{k}{r}}$. So we get the first result. 
    \item Since $h=h_1\cdots h_r$ we have $\vsg(h)=\prod_{i=1}^r\vsg(h_i)$. Now by \thmref{thm:irr-facts-into-irr-permutes} we have $\vsg(h_i)\in E$ for all $i\in[r]$ and they are conjugates of each other over $\bbF_q$. Therefore $\vsg(h)=N_{E/\bbF_q}(\vsg(h_i))$ for any $i\in[r]$. Therefore $$N_{K/\bbF_q}(\vsg(h_i))=\lt(N_{E/\bbF_q}(h_i)\rt)^{\nicefrac{k}{r}}=\vsg(h)^{\nicefrac{k}{r}}\quad \forall\ i\in[r]$$On the other hand we have $\vrh(h)=\sum_{i=1}^r\vsg(h_i)$ for any $i\in[r]$. Therefore $\vrh(h)=\tr_{E/\bbF_q}$. Thus $$\tr_{K/\bbF_q}(\vsg(h_i)) =\frac{k}{r}\tr_{E/\bbF_q}(\vrh(h_i))=\frac{k}{r}\vrh(r)$$So together we have $$\xi^{(k)}(h_i)= \psi^{(k)}\circ\vsg(h_i)\cdot \chi^{(k)}\circ\vrh(h)=\psi\lt(N_{K/\bbF_q}(\vsg(h_i))\rt)\cdot \chi\lt(\tr_{K/\bbF_q}(\vrh(h_i))\rt)=\psi^{\nicefrac{k}r}(\vsg(h))\cdot \chi^{\nicefrac{k}r}(\vsg(h))=\xi^{\nicefrac{k}{r}}(h)$$So we have the theorem. 
  \end{enumerate}
\end{proof}
Now we can write $L^{(k)}(s,\xi^{(k)})$ and $\ovL^{(k)}(z,\xi^{(k)})$ in terms of $L$-functions on the field $\bbF_q$ using $\xi$ and factorize the lifted $L$-function as product of $L$-function on the base field. 
\begin{Theorem}{Factorisation of the Lifted $L$-function}{thm:lifted-L-factorisation}
  Let $K=\bbF_{q^k}$ be a finite extension of $\bbF_q$ and let $\zeta_k=e^{2\pi i/k}$ be a primitive $k$-th root of unity. Then
  $$L^{(k)}\lt(s,\xi^{(k)}\rt)=\prod_{t=1}^{k}L\lt(s-\frac{2\pi it}{k\log q},\xi\rt).$$
\end{Theorem}
\begin{proof}
  By \LEulerProdthm we have $$L^{(k)}\lt(s,\xi^{(k)}\rt)=\prod_{l(X)\in\irr(\Phi^{(k)})}\lt(1-\xi^{(k)}\mfn^{-s}{(k)}(l)\rt)^{-1}$$Since $h$ is an irreducible monic polynomial in $\bbF_q[X]$ of degree $d$ it splits intro $r$ polynomials over $K$ where $r=\gcd(d,k)$ i.e. $h=h_1\cdot h_r$. Now by \IrrSplitslem $$\xi^{(k)}(h_i)\mfn^{-s}{(k)}(h_i)=\lt(\xi(h)\mfn^{-s}(h)\rt)^{\nicefrac{k}{r}}$$On the other hand, every monic irreducible polynomial $l(X)\in k[X]$ is the factor of an unique monic irreducible $h(X)\in\bbF_q[X]$. Therefore $$L^{(k)}\lt(s,\xi^{(k)}\rt)=\prod_{h\in\irr(\Phi)}\lt(1-\lt(\xi(h)\mfn^{-s}(h)\rt)^{\frac{k}{\gcd(k,\deg(h))}}\rt)^{-\gcd(k,\deg(h))}$$Then by \RootsUnityProdlem we have 
  \begin{align*}
    L^{(k)}\lt(s,\xi^{(k)}\rt) & = \prod_{h\in\irr(\Phi)}\lt(1-\lt(\xi(h)\mfn^{-s}(h)\rt)^{\frac{k}{\gcd(k,\deg(h))}}\rt)^{-\gcd(k,\deg(h))}\\ 
    & = \prod_{h\in\irr(\Phi)}\prod_{t=1}^k\lt(1-e^{2\pi i\frac{t\cdot \deg(h)}{k}}\xi(h)\mfn^{-s}(h)\rt)^{-1}\\ 
    & = \prod_{t=1}^k \prod_{h\in\irr(\Phi)}\lt(1-q^{2\pi i\frac{t\cdot \deg(h)}k \log_qe}\xi(h)\mfn^{-s}(h)\rt)^{-1}\\ 
    & = \prod_{t=1}^k \prod_{h\in\irr(\Phi)}\lt(1-\xi(h)\mfn^{-\lt(s-2\pi i\frac{t}{k\log q}\rt)}(h)\rt)^{-1}\\ 
    & = \prod_{t=0}^{k-1}L\lt(s-\frac{2\pi it}{k\log q},\xi\rt)
  \end{align*}So we have the theorem. 
\end{proof}
Now by \LXiFinitethm, $\Lxi$ is a polynomial. So there exist complex numbers $w_1,\dots, w_l\in\bbC$ such that $\Lxi=(1-w_1u)\cdots (1-w_lu)$ for some $l\in\bbN$ where $u=q^{-s}$. By \LiftedLFactthm, $\ovLxil{k}$ or $\Lxil{k}$ is also a polynomial (a product of finitely many polynomials). We will prove that the reciprocal roots of $\Lxil{k}$ are $w_i^k$ for all $i\in[l]$. 
\begin{corolary}{Reciprocal Roots of the Lifted $L$-function}{cor:lifted-L-roots}
  If $\Lxi=(1-w_1u)\cdots (1-w_lu)$ with $u=q^{-s}$ then $$\Lxil{k}=(1-w_1^k\cdot u_k)\cdots (1-w_l^k\cdot u_k)$$where $u_k=q^{-ks}$.
\end{corolary}
\begin{proof}
  By \LiftedLFactthm we have $$\Lxil{k}=\prod_{t=1}^k\Lxis{s-\frac{2\pi it}{k\log q}}$$ Now for any $t\in[k]$ we have $q^{-\lt(s-2\pi i\frac{t}{k\log q}\rt)}=e^{2\pi i \frac{t}{k}}\cdot q^{-s}=e^{2\pi i\frac{t}{k}}\cdot u$. So we get
  \begin{align*}
    \Lxil{k} & = \prod_{t=1}^k \Lxis{s-\frac{2\pi it}{k\log q}}\\
    & = \prod_{t=1}^k \prod_{j=1}^l \lt(1-w_je^{2\pi i\frac{t}{k}}u\rt)\\
    & = \prod_{j=1}^l \prod_{t=1}^k \lt(1-w_je^{2\pi i\frac{t}{k}}u\rt)\\
    & = \prod_{j=1}^l \lt(1-w_j^ku^k\rt)\\
    & = \prod_{j=1}^l \lt(1-w_j^k\cdot u_k\rt)
  \end{align*}
  So we have the result. 
\end{proof}

So now we can give an analogue of \LLeadCoeffthm for the lifted sum and give a relation with $\Wmixl{k}{f}{g}$ with the coefficients of $\Lxil{k}$. 
\begin{corolary}{Character Sum via Reciprocal Roots}{cor:charsum-reciprocal-roots}
  Suppose either $\psi\in\Mq$ is a nontrivial multiplicative character of exponent $d$ and $Y^d-f(X)$ is absolutely irreducible. If $\chi\in\Xq$ is a additive character either non-trivial or trivial and $f(X)=1$ then $$W^{(k)}\lt(\psi^{(k)},\chi^{(k)};f,g\rt)=-\sum_{i=1}^{n+m-1}w_i^k.$$
\end{corolary}
\begin{proof}
  If we apply \LLeadCoeffthm to $\bbF_{q^k}$ instead of $\bbF_q$ then using \LiftedLRootscor we have $$\Lxil{k}=1+c_{k,1}\cdot u_k+\cdots +c_{k,n+m-1}\cdot u_k^{n+m-1}=\prod_{j=1}^{n+m-1}(1-w_j^k\cdot u_k)$$Thereore $c_{k,1}=\Wmixl{k}{f}{g}$. And on the other hand we have that $c_{k,1}=-\sum_{j=1}^{n+m-1}w_j^k$. 
\end{proof}
\subsection{Alternate Proof of Davenport-Hasse Relation}